\subsection{ms-usuarios --- Tenants, periodos, usuarios y perfiles}
\label{sec:ms-usuarios}
\textit{Responsable: Benjamín (ACA-01) y Edinson (ACA-02). Repositorio: \url{https://github.com/nubeduca/nubeduca-ms-usuarios}.}

ms-usuarios es la base del resto: guarda las instituciones (tenants), sus periodos académicos y el perfil de cada usuario. Los tres tenants (UTEC, UPC y PUCP) los crea el seed y emite un \texttt{tenant.created} por cada uno. El administrador crea periodos en \texttt{PLANNED} y los inicia (\texttt{ACTIVE}); al iniciarlos se emite \texttt{term.started}, y ms-cursos solo acepta cursos en periodos activos. El nombre de un periodo es único en su tenant gracias a un ítem de reserva \texttt{TERMNAME\#} que se escribe en la misma transacción.

El perfil se crea en el primer login: el frontend llama a \texttt{POST /users/me} justo después de autenticarse. El tenant, el rol y el \texttt{userId} salen del token de Cognito, nunca del cuerpo; el correo y el nombre salen de los \textit{claims}, y lo que falte, junto con el código institucional, sale del padrón \texttt{ROSTER\#<email>} que la institución carga de antemano. La llamada es idempotente: la primera vez responde 201 y emite \texttt{user.provisioned}; las siguientes devuelven el perfil existente con 200 y sin evento. El evento no lleva nombre ni correo, solo el id, el rol y el código: ms-cursos lo usa para traducir el código institucional del CSV de matrícula a un \texttt{userId}, y ms-notificaciones para crear las preferencias por defecto. Un ítem \texttt{USERCODE\#} garantiza que el código no se repita dentro del tenant.

Todos los roles leen su tenant, los periodos y su propio perfil, y pueden cambiar solo su idioma y zona horaria (\texttt{PATCH /users/me}). Docentes y administradores pueden listar y buscar usuarios del tenant, por rol o por código, por ejemplo para armar una matrícula; solo el administrador crea e inicia periodos. Como todo se guarda en la partición de su \texttt{tenant\_id}, un usuario de otra institución recibe 404 aunque conozca el id. ms-usuarios no consume eventos.

\subsubsection{Endpoints}
\begin{table}[htbp]
\centering\footnotesize
\begin{tabularx}{\textwidth}{lXll}
\toprule
Método & Ruta & Roles & Emite \\
\midrule
GET    & \texttt{/health} & pública & --- \\
GET    & \texttt{/tenant} & todos & --- \\
GET    & \texttt{/terms} & todos & --- \\
POST   & \texttt{/terms} & ADMIN & --- \\
POST   & \texttt{/terms/\{termId\}/start} & ADMIN & \texttt{term.started} \\
POST   & \texttt{/users/me} & todos & \texttt{user.provisioned} \\
GET    & \texttt{/users/me} & todos & --- \\
PATCH  & \texttt{/users/me} & todos & --- \\
GET    & \texttt{/users} & TEACHER, ADMIN & --- \\
GET    & \texttt{/users/\{userId\}} & TEACHER, ADMIN & --- \\
\bottomrule
\end{tabularx}
\caption{Endpoints de ms-usuarios (\texttt{apis/ms-usuarios.openapi.yaml} de \texttt{nubeduca-contratos}); \textit{todos} = STUDENT, TEACHER y ADMIN.}
\end{table}

\subsubsection{Eventos}
\textbf{Publica:}
\begin{itemize}
  \item \texttt{lms.users.tenant.created.v1}
  \item \texttt{lms.users.term.started.v1}
  \item \texttt{lms.users.user.provisioned.v1}
\end{itemize}

\textbf{Consume:}
\textit{Ninguno.}


\subsubsection{Estructura JSON de la tabla \texttt{<stage>-t\_usuarios}}
Clave de partición \texttt{tenant\_id}, clave de ordenamiento \texttt{sk}.

\begin{table}[htbp]
\centering\small
\begin{tabularx}{\textwidth}{>{\raggedright\arraybackslash}p{0.42\textwidth}X}
\toprule
\texttt{sk} & Qué es \\
\midrule
\texttt{TENANT} & Institución \\
\texttt{TERM\#<termId>} & Periodo académico \\
\texttt{TERMNAME\#<name>} & Reserva del nombre del periodo en el tenant \\
\texttt{USER\#<userId>} & Perfil \\
\texttt{USERCODE\#<código>} & Código institucional $\rightarrow$ usuario (único en el tenant) \\
\texttt{ROSTER\#<email>} & Padrón: lo que la institución ya sabe de ese correo, para el primer login \\
\texttt{EVENTO\#<eventId>} & Outbox (lo lee el Stream) \\
\bottomrule
\end{tabularx}
\caption{Prefijos de \texttt{sk} en la tabla de ms-usuarios.}
\end{table}

\begin{lstlisting}
{
  "tenant_id": "utec",
  "sk": "TENANT",
  "entidad": "TENANT",
  "name": "Universidad de Ingeniería y Tecnología",
  "timeZone": "America/Lima",
  "locale": "es-PE",
  "version": 1,
  "createdAt": "2026-11-13T15:00:00Z"
}
\end{lstlisting}

\begin{lstlisting}
{
  "tenant_id": "utec",
  "sk": "TERM#trm_01JCP0KWJDXF9GJ9YNTK2H84A9",
  "entidad": "TERM",
  "termId": "trm_01JCP0KWJDXF9GJ9YNTK2H84A9",
  "name": "2026-2",
  "startsOn": "2026-08-17",
  "endsOn": "2026-12-12",
  "status": "ACTIVE",
  "version": 2,
  "createdBy": "16248d98-4e7a-41b9-a08b-4929177d2e2d",
  "createdAt": "2026-11-13T15:00:00Z",
  "startedBy": "16248d98-4e7a-41b9-a08b-4929177d2e2d",
  "startedAt": "2026-11-13T15:05:00Z"
}
\end{lstlisting}

\begin{lstlisting}
{
  "tenant_id": "utec",
  "sk": "USER#11111111-1111-4111-8111-111111111111",
  "entidad": "USER",
  "userId": "11111111-1111-4111-8111-111111111111",
  "role": "STUDENT",
  "institutionalCode": "202310457",
  "givenName": "Ana",
  "familyName": "Torres",
  "email": "alumno.demo@utec.edu.pe",
  "locale": "es-PE",
  "timeZone": "America/Lima",
  "version": 1,
  "createdAt": "2026-11-10T15:00:00Z"
}
\end{lstlisting}

\begin{lstlisting}
{
  "tenant_id": "utec",
  "sk": "ROSTER#alumno.demo@utec.edu.pe",
  "entidad": "ROSTER",
  "institutionalCode": "202310457",
  "givenName": "Ana",
  "familyName": "Torres",
  "locale": "es-PE",
  "timeZone": "America/Lima"
}
\end{lstlisting}

Formato de los ítems documentado en el README de \texttt{nubeduca-ms-usuarios}.

\subsubsection{Datos de prueba}
El script \texttt{scripts/poblar\_datos.py} crea los tres tenants (UTEC, UPC y PUCP) y publica un \texttt{tenant.created} por cada uno; si el tenant ya existe, no lo publica de nuevo. Después carga 30\,000 perfiles: UTEC 15\,000, UPC 10\,000 y PUCP 5\,000, con 90\,\% de alumnos, 9\,\% de docentes y 1\,\% de administradores, cada uno con su código institucional. Los perfiles no escriben ítems \texttt{EVENTO\#}, para no inundar el bus, y sus identificadores son fijos, así que el script es idempotente. Incluye los usuarios demo del frontend y un periodo académico activo de seed, el mismo que usa ms-cursos.

Carga en la tabla desplegada \texttt{dev-t\_usuarios} (22 de noviembre de 2026):

\begin{lstlisting}
$ python scripts/poblar_datos.py --stage dev
utec: tenant CREADO
upc: tenant CREADO
pucp: tenant CREADO
utec: 15000 usuarios
upc: 10000 usuarios
pucp: 5000 usuarios
Total: 30000 usuarios en dev-t_usuarios
\end{lstlisting}

El conteo de los perfiles (\texttt{USER\#}) por tenant en la tabla desplegada, con \texttt{aws dynamodb query ... --select COUNT}, confirma UTEC 15\,000, UPC 10\,000 y PUCP 5\,000 (figura~\ref{fig:ms-usuarios-datos}).

\evidencia{ms-usuarios-datos.png}{ms-usuarios: conteo de perfiles por tenant en \texttt{dev-t\_usuarios}.}{fig:ms-usuarios-datos}

\subsubsection{Resultados y evidencias}
\evidencia{ms-usuarios-evidencia.png}{ms-usuarios desplegado en AWS: \texttt{GET /users?role=TEACHER\&limit=3} como administrador de UTEC devuelve 200 con docentes del seed.}{fig:ms-usuarios}
